\documentclass[twoside,12pt]{book}

\usepackage[utf8]{inputenc}
\usepackage[T1]{fontenc}
\usepackage{mathpazo}
\usepackage{amsmath,amssymb,amsthm}
\usepackage{graphicx}
\usepackage{listings}
\usepackage{xcolor}
\usepackage{enumitem}
\usepackage{tikz}
\usepackage{algorithm}
\usepackage{algpseudocode}
\usepackage{booktabs}
\usepackage{geometry}
\usepackage{fancyhdr}
\usepackage{titlesec}
\usepackage[colorlinks=true, linkcolor=blue, citecolor=blue, urlcolor=blue, plainpages=false, pdfpagelabels]{hyperref}

\geometry{
  papersize={6in,9in},
  margin=0.75in,
  top=0.75in,bottom=0.75in,
  inner=0.85in,outer=0.65in,
  headheight=14pt
}

\pagestyle{fancy}
\fancyhf{}
\fancyhead[LE]{\small\textit{1. Introduction}}
\fancyhead[RO]{\small\textit{Randomized Algorithms}}
\fancyfoot[C]{\thepage}
\renewcommand{\headrulewidth}{0.4pt}

\definecolor{codegray}{rgb}{0.45,0.45,0.45}
\lstdefinestyle{cppstyle}{
  basicstyle=\ttfamily\scriptsize,
  keywordstyle=\bfseries,
  commentstyle=\itshape\color{codegray},
  numbers=left,
  numberstyle=\tiny\color{codegray},
  numbersep=8pt,
  frame=single,
  framerule=0.4pt,
  rulecolor=\color{codegray},
  backgroundcolor=\color{gray!5},
  breaklines=true,
  tabsize=4,
  showstringspaces=false,
  language=C++,
  morekeywords={nullptr,size_t,auto,const,constexpr,
    unordered_map,vector,pair,mt19937,
    uniform_int_distribution,INT_MAX}
}
\lstset{style=cppstyle}

\theoremstyle{plain}
\newtheorem{theorem}{Theorem}[chapter]
\newtheorem{lemma}[theorem]{Lemma}
\newtheorem{corollary}[theorem]{Corollary}
\newtheorem{proposition}[theorem]{Proposition}
\theoremstyle{definition}
\newtheorem{definition}[theorem]{Definition}
\newtheorem{example}[theorem]{Example}
\newtheorem{remark}[theorem]{Remark}
\theoremstyle{remark}
\newtheorem{exercise}[theorem]{Exercise}
\newtheorem{problem}[theorem]{Problem}
\newenvironment{cpplisting}{\begin{center}\small}{\end{center}}

\newcommand{\Exp}{\operatorname{\mathbb{E}}}
\newcommand{\Var}{\operatorname{Var}}
\newcommand{\Cov}{\operatorname{Cov}}

\begin{document}

\chapter*{1 Introduction}
\addcontentsline{toc}{chapter}{1 Introduction}
\markboth{1 Introduction}{Randomized Algorithms}
\label{ch:introduction}

Algorithms are the engines of computation---systematic procedures that
transform inputs into outputs through a sequence of well-defined steps.
For decades, the gold standard in algorithm design has been
\emph{determinism}: given the same input, a deterministic algorithm always
produces the same output and takes the same time. Yet there exists a large
class of problems for which no efficient deterministic algorithm is known,
and for which the introduction of randomness yields striking improvements in
simplicity, speed, or both. This book is devoted to the theory and practice
of \emph{randomized algorithms}---algorithms that make use of random choices
during their execution.

We begin this chapter by examining why deterministic algorithms, despite
their elegance and predictability, face fundamental limitations. This
motivates the study of randomization as a design principle. We then
discuss random number generation in practice, since every randomized
algorithm ultimately relies on a source of randomness. Next, we gather
the elementary tools from probability theory that will recur throughout
the book: conditional probability, linearity of expectation, and tail
inequalities. Finally, we introduce the main ideas behind randomized
algorithms, covering the two principal paradigms---Las Vegas and Monte
Carlo methods---and illustrating them with concrete examples.

%==========================================================================
\section{Limitations of Deterministic Algorithms}
\label{sec:det-limits}
%==========================================================================

To appreciate the power of randomization, it is essential to understand
where deterministic methods fall short. By ``deterministic algorithm'' we
mean an algorithm whose behavior is completely determined by its input: no
coin flips, no random choices. The running time of a deterministic algorithm
is a function of the input alone.

\subsection{Worst-case complexity is not the full picture}

The traditional framework for analyzing algorithms is \emph{worst-case
complexity}: we measure the running time on the most adversarial input of
size~$n$. This framework has been enormously productive, yielding tight
bounds for many problems (e.g.\ sorting in $O(n \log n)$ time, matrix
multiplication via Strassen's algorithm). However, worst-case analysis can
be overly pessimistic. A deterministic algorithm may be designed to protect
against a pathological input that rarely---or never---occurs in practice.
The resources devoted to handling the worst case may therefore be wasted on
typical inputs.

\subsubsection{Any order versus random order}

A subtle but important distinction concerns the order in which input
elements are presented to an algorithm. This becomes particularly
significant for incremental algorithms that process elements one by one.

\textbf{Any order (adversarial).} When a deterministic incremental algorithm
processes elements~\emph{in any order}, an adversary who knows the algorithm
can choose the order that maximizes the running time or solution cost. For
example, inserting points into a convex hull in worst-case order causes
every insertion to restructure the hull, yielding $O(n^2)$ time. Similarly,
adding elements to a hash table in adversarial order can cause all elements
to collide.

\textbf{Random order.} If the same elements are processed in a uniformly
random permutation, the expected behavior is often dramatically better.
Randomized incremental algorithms exploit this idea: randomizing the order
in which elements are processed turns worst-case inputs into average-case
behavior. For convex hulls, a random insertion order yields $O(n \log n)$
expected time (Section~10.2). For linear programming, the randomized
incremental algorithm of Seidel runs in expected $O(d! n)$ time, which is
linear in~$n$ for fixed dimension~$d$ (Section~10.9).

This principle---\emph{randomizing the order of processing eliminates the
adversary's power to choose a bad order}---is one of the most practical
benefits of randomization in algorithm design. It requires no change to the
input itself, only to the algorithm's scheduling.

\subsection{Problems with no known efficient deterministic solution}

There are problems for which no polynomial-time deterministic algorithm is
known, despite intensive research over many decades. Three categories
illustrate this:

\begin{enumerate}
\item \textbf{NP-hard problems.} The class NP encompasses decision problems
whose YES instances admit short certificates verifiable in polynomial time.
NP-hard problems are at least as hard as every problem in NP. No
polynomial-time algorithm for any NP-hard problem is known, and the
existence of such an algorithm would imply $\mathbf{P} = \mathbf{NP}$.
Famous NP-hard problems include the traveling salesman problem, graph
coloring, and satisfiability (SAT). While no algorithm---deterministic or
randomized---is known to solve these efficiently in general, randomization
yields algorithms that perform well on average or that produce approximate
solutions with provable guarantees (Chapter~12).

\item \textbf{Problems requiring global information.} Some problems seem
inherently difficult for deterministic algorithms because a deterministic
procedure cannot ``explore'' the solution space without being trapped by
adversarial structure. Randomized algorithms can avoid such traps by
exploring the space in an unpredictable fashion.

\item \textbf{Adversarial inputs.} A deterministic algorithm, once its
strategy is known, can be defeated by an adversary who constructs a
worst-case input tailored to exploit the algorithm's predictable behavior.
Randomness removes this vulnerability: an adversary cannot predict the
algorithm's random choices.
\end{enumerate}

\subsection{The adversarial argument}

The idea that an adversary can defeat a deterministic algorithm is
formalized in the following observation. Let $A$ be a deterministic
algorithm for a problem~$\Pi$. Suppose there exists an input $x$ on which
$A$ takes super-polynomial time or produces an incorrect answer. Then the
adversary who knows $A$ can always present $x$ as the input. The
deterministic algorithm has no recourse---it will always behave badly on
$x$.

A randomized algorithm $R$, by contrast, makes random choices during its
execution. Even if the adversary knows the code of $R$ and the input $x$,
the adversary cannot predict the outcome of $R$'s random coins. The best
the adversary can do is ensure that $R$ performs badly on \emph{some} set
of random outcomes; but if $R$ is well designed, the set of ``bad''
outcomes has small probability.

This is the fundamental insight behind the study of randomized algorithms:
\emph{randomness provides immunity against adversarial worst-case inputs.}

\subsection{A concrete example: finding a median}

Consider the problem of finding the median of $n$ numbers. The classical
deterministic algorithm of Blum, Floyd, Pratt, Rivest, and Tarjan solves
this in $O(n)$ time, but the constant hidden in the $O$-notation is
large---the algorithm is complex and rarely used in practice. The
``textbook'' deterministic algorithm based on sorting takes $O(n \log n)$
time.

A simple randomized algorithm achieves $O(n)$ expected time with
remarkable ease: choose a random element as a pivot, partition the array,
and recurse into the side containing the median. The analysis (see
Chapter~4) shows that the expected number of comparisons is at most $4n$.
This algorithm is not only simpler than its deterministic counterpart but
also faster in practice, despite being randomized. (A deterministic
linear-time median algorithm with smaller constants was later discovered,
but the randomized version remains a compelling illustration of the power
of simplicity.)

%==========================================================================
\section{Random Number Generation}
\label{sec:rng}
%==========================================================================

Every randomized algorithm ultimately depends on a source of randomness.
In theoretical analysis we model randomness as an infinite sequence of
independent, uniformly distributed random bits. In practice, we need
algorithms---\emph{pseudorandom number generators} (PRNGs)---that produce
sequences that are ``random enough'' for the application at hand.

\subsection{Pseudorandom number generators}

A PRNG is a deterministic procedure that, given an initial value called a
\emph{seed}, produces a sequence of numbers that appears statistically
random. The sequence is entirely determined by the seed; two runs with the
same seed produce identical sequences.

Modern PRNGs are designed to satisfy two desiderata:
\begin{enumerate}
\item \textbf{Statistical quality:} The output should be indistinguishable
from truly random bits by efficient statistical tests.
\item \textbf{Efficiency:} The generation of each pseudorandom number
should take $O(1)$ time.
\end{enumerate}

The most widely used PRNG in practice is the \emph{Mersenne Twister}
(mt19937), which produces 32-bit pseudorandom integers with a period of
$2^{19937}-1$ and passes most standard statistical tests. C++11 provides
\texttt{std::mt19937} and related distributions via the \texttt{<random>}
header.

\begin{lstlisting}[
  caption={Using the C++ <random> library.},
  label={lst:cpprandom}]
#include <random>
#include <chrono>
#include <iostream>

int main() {
  // Seed from system clock
  unsigned seed = static_cast<unsigned>(
    std::chrono::steady_clock::now()
      .time_since_epoch().count());
  std::mt19937 eng(seed);

  // Uniform integer in [1, 100]
  std::uniform_int_distribution<int> d1(1,100);
  std::cout << "Dice roll: " << d1(eng) << "\n";

  // Uniform real in [0.0, 1.0)
  std::uniform_real_distribution<double> d2(0.0,1.0);
  std::cout << "Uniform:   " << d2(eng) << "\n";

  // Normal distribution
  std::normal_distribution<double> d3(0.0,1.0);
  std::cout << "Gaussian:  " << d3(eng) << "\n";
}
\end{lstlisting}

\subsection{Seeding and reproducibility}

The choice of seed is critical for reproducibility. In scientific
simulations and debugging, it is common to fix the seed so that results can
be reproduced exactly. In production systems, the seed is typically drawn
from a source of entropy (e.g., the system clock, hardware noise, or
\texttt{/dev/urandom} on Unix systems).

Using the same seed with the same PRNG always yields the same sequence.
This property is invaluable for debugging: a bug that manifests on a
particular input can be reproduced by using the same seed.

\subsection{Cryptographically secure PRNGs}

For cryptographic applications, standard PRNGs are insufficient: an attacker
who observes a portion of the output should not be able to predict future
output. Cryptographically secure PRNGs (CSPRNGs) are designed to resist such
attacks. On most systems, \texttt{/dev/urandom} (Unix) or
\texttt{BCryptGenRandom} (Windows) provide CSPRNG output.

\subsection{Weak randomness versus strong randomness}

Not all applications require the same quality of randomness. This
distinction leads to two categories of random sequences.

\subsubsection{Weak random sequences}

A \emph{weak random sequence} is one generated by a simple PRNG with modest
statistical guarantees. Examples include linear congruential generators
(LCGs) and the standard C \texttt{rand()} function. Weak PRNGs are fast and
require little state, but their output has detectable correlations: LCGs
suffer from lattice structure in higher dimensions, and short periods cause
repeat patterns in long runs.

Weak randomness suffices for:
\begin{itemize}
\item \textbf{Monte Carlo integration.} Estimating $\pi$ by sampling points
in a square (Section~1.4.7) converges at $O(1/\sqrt{n})$ regardless of mild
correlations in the samples.
\item \textbf{Randomized incremental algorithms.} Randomizing the insertion
order for convex hulls or Delaunay triangulations (Chapter~10) only requires
that each permutation be approximately equally likely; a weak PRNG suffices.
\item \textbf{Approximate counting.} The Karp--Luby scheme for DNF counting
(Chapter~12) relies on the expectation of indicator variables; the analysis
holds as long as samples are pairwise independent.
\end{itemize}

However, weak randomness can silently invalidate results in applications
that demand uniformity over large state spaces. The RANDU generator (a
notorious LCG from the 1960s) produced triples $(x, W(x), W(W(x)))$ that
all lay on 15 planes in 3D space, corrupting 3D Monte Carlo simulations.

\subsubsection{Strong random sequences}

A \emph{strong random sequence} is one that is statistically
indistinguishable from truly random bits by any polynomial-time test.
The Mersenne Twister (mt19937) with its $2^{19937}-1$ period qualifies for
most non-cryptographic uses. Cryptographically secure PRNGs (CSPRNGs)
provide the strongest guarantee: no polynomial-time adversary can predict
future bits given access to previous output.

Strong randomness is required for:
\begin{itemize}
\item \textbf{Cryptographic protocols.} Key generation, zero-knowledge
proofs, and commitment schemes (Section~15.4) require unpredictability
against an active adversary. A weak PRNG-derived RSA key can be factored.
\item \textbf{Derandomization.} The method of conditional probabilities
(Chapter~6) and pairwise independence (Chapter~4) require true random
bits or provably strong generators.
\item \textbf{Long-running simulations.} Simulations of $10^{12}$ steps
exhaust the period of a 32-bit LCG ($\approx 4\times 10^9$). The
Mersenne Twister's period of $2^{19937}$ is effectively inexhaustible.
\end{itemize}

\subsubsection{Effect on algorithmic guarantees}

A randomized algorithm's theoretical guarantees are proven under the
assumption of true randomness (independent, uniformly distributed bits).
When implemented with a PRNG, the empirical behavior may differ:

\begin{itemize}
\item \textbf{Las Vegas algorithms.} The running time distribution may
shift slightly, but correctness is unaffected: the algorithm always
produces the right answer regardless of the random choices.
\item \textbf{Monte Carlo algorithms.} The error probability may increase
if the PRNG introduces correlations. For example, using \texttt{rand()}
with the Miller--Rabin primality test can misclassify Carmichael numbers
more frequently than the theoretical bound $4^{-k}$.
\item \textbf{Probabilistic method constructions.} The non-constructive
existence proofs (Chapter~6) are unaffected by the PRNG quality, but any
explicit construction derived from derandomization requires strong
randomness.
\end{itemize}

\begin{remark}
In practice, the Mersenne Twister (mt19937) provides sufficient randomness
for nearly all non-cryptographic randomized algorithms discussed in this
book. For cryptographic applications (Chapter~15), always use a CSPRNG from
the operating system's entropy pool.
\end{remark}

\subsection{True randomness versus pseudorandomness}

No sequence produced by a deterministic algorithm is truly random. Every
PRNG is a finite-state machine: given the same seed, it always produces the
same output. The sequence only \emph{appears} random to computational tests.

\subsubsection{Sources of true randomness}

True randomness exists in physical processes:
\begin{itemize}
\item \textbf{Quantum phenomena.} Radioactive decay, photon detection, and
shot noise in semiconductor junctions are fundamentally unpredictable.
Hardware random number generators (HRNGs) exploit these effects.
\item \textbf{Atmospheric noise.} Radio static from lightning strikes
provides a continuous source of entropy.
\item \textbf{Human interaction.} Mouse movements, keyboard timings, and
disk seek latencies are unpredictable enough for seeding PRNGs.
\end{itemize}

Modern operating systems maintain entropy pools fed by hardware interrupts.
On Linux, \texttt{/dev/random} blocks when entropy is low;
\texttt{/dev/urandom} does not but uses a CSPRNG to expand the entropy.

\subsubsection{Practical implications}

For most applications, a well-seeded PRNG is indistinguishable from true
randomness. However, three situations require genuine entropy:
\begin{enumerate}
\item \textbf{Cryptographic key generation.} A key generated by a PRNG with
a predictable seed can be recovered by an adversary.
\item \textbf{Scientific reproducibility.} True randomness cannot reproduce
a simulation; PRNGs with saved seeds enable exact reproducibility.
\item \textbf{Algorithmic fairness.} Lotteries, elections, and audit
selection must resist manipulation. Physical entropy sources (or
public-coin protocols) provide verifiable randomness.
\end{enumerate}

\subsubsection{Buffon's needle: the first Monte Carlo method}

The earliest documented use of randomness to compute a mathematical
constant predates electronic computers by two centuries.

\begin{example}[Buffon's Needle, 1777]
A needle of length~$\ell$ is dropped onto a floor made of parallel wooden
strips of width~$d$ (with $\ell \le d$). What is the probability that the
needle crosses a crack between strips?

Let $y \in [0, d/2]$ be the distance from the needle's midpoint to the
nearest crack, and $\theta \in [0, \pi/2]$ be the acute angle between the
needle and the crack direction. The needle crosses a crack iff
$y \le (\ell/2) \sin\theta$. Assuming $y$ and $\theta$ are independent and
uniformly distributed, the probability is
\[
  p = \frac{2}{\pi} \cdot \frac{\ell}{d}.
\]

By repeating the experiment $n$ times and counting the number $m$ of
crossings, we estimate $\pi$ as
\[
  \pi \approx \frac{2\ell n}{d m}.
\]
\end{example}

Buffon's needle is historically significant for three reasons:
\begin{enumerate}
\item It is the first recorded \emph{Monte Carlo method}---using random
samples to compute a deterministic quantity.
\item It demonstrates that randomness can evaluate integrals without
analytic solution, a principle now central to computational physics,
finance, and Bayesian statistics.
\item The convergence rate $O(1/\sqrt{n})$ is characteristic of all Monte
Carlo methods and sets a fundamental limit on their accuracy.
\end{enumerate}

The needle problem also illustrates a recurring theme: a randomized
algorithm's guarantees rely on the quality of the random source. If the
dropped positions are biased (e.g.\ the dropper unconsciously avoids
cracks), the estimate of $\pi$ becomes systematically wrong. Similarly, a
PRNG with low-period correlations can introduce bias in Monte Carlo
simulations that is invisible to simple statistical tests.

\subsection{Random bits vs.\ random numbers}

Most randomized algorithms require random integers or reals, not raw random
bits. From a theoretical standpoint, this distinction is immaterial: a
random $k$-bit integer can be obtained by sampling $k$ independent random
bits. In practice, PRNGs are often optimized for generating integers in a
given range, which is more efficient than generating raw bits and reducing
modulo.

\begin{remark}
Throughout this book, we assume the availability of a source of independent,
uniformly distributed random bits. When analyzing randomized algorithms, we
study the probability distribution over all possible random outcomes. The
quality of the PRNG used in an implementation affects the empirical
performance but not the theoretical guarantees (which hold for truly random
coins).
\end{remark}

%==========================================================================
\section{Preliminary Mathematics}
\label{sec:prelim-math}
%==========================================================================

The analysis of randomized algorithms draws heavily on elementary
probability theory. In this section we review the tools that will be used
throughout the book. More advanced material (Chernoff bounds, martingales,
the probabilistic method) is introduced in later chapters as needed.

\subsection{Sample spaces and events}

A \emph{probability space} is a triple $(\Omega, \mathcal{F}, \Pr)$,
where $\Omega$ is a set (the \emph{sample space}), $\mathcal{F}$ is a
collection of subsets of $\Omega$ (the \emph{events}), and $\Pr$ is a
probability measure satisfying $\Pr[\Omega]=1$ and countable additivity.

For finite sample spaces, we typically take $\mathcal{F}$ to be all subsets
of $\Omega$ and assign a probability $\Pr[\omega]$ to each outcome
$\omega \in \Omega$ such that $\sum_{\omega \in \Omega} \Pr[\omega] = 1$.
For a uniform distribution, every outcome has probability $1/|\Omega|$.

\subsection{Independence and conditional probability}

Two events $\mathcal{E}_1$ and $\mathcal{E}_2$ are said to be
\emph{independent} if the probability that they both occur is given by
\begin{equation*}
  \Pr[\mathcal{E}_1 \cap \mathcal{E}_2]
  = \Pr[\mathcal{E}_1] \times \Pr[\mathcal{E}_2].
  \tag{1.3}
\end{equation*}
In the more general case where $\mathcal{E}_1$ and $\mathcal{E}_2$ are not
necessarily independent,
\begin{equation*}
  \Pr[\mathcal{E}_1 \cap \mathcal{E}_2]
  = \Pr[\mathcal{E}_1] \times \Pr[\mathcal{E}_2 \mid \mathcal{E}_1],
  \tag{1.4}
\end{equation*}
where $\Pr[\mathcal{E}_1 \mid \mathcal{E}_2]$ denotes the conditional
probability of $\mathcal{E}_1$ given $\mathcal{E}_2$. The chain rule for
probabilities extends this to $k$ events:
\begin{equation*}
  \Pr\!\Bigl[\bigcap_{i=1}^{k} \mathcal{E}_i\Bigr]
  = \Pr[\mathcal{E}_1]
    \times \Pr[\mathcal{E}_2 \mid \mathcal{E}_1]
    \times \Pr[\mathcal{E}_3 \mid \mathcal{E}_1 \cap \mathcal{E}_2]
    \cdots
    \Pr\!\Bigl[\mathcal{E}_k \mid \bigcap_{i=1}^{k-1} \mathcal{E}_i\Bigr].
  \tag{1.5}
\end{equation*}

We used this chain rule in deriving the success probability of Karger's
min-cut algorithm (Eq.~1.1).

\subsection{Expectation and linearity}

The \emph{expected value} (or \emph{mean}) of a discrete random variable
$X$ taking values $x_1, x_2, \ldots$ with probabilities $p_1, p_2,
\ldots$ is
\[
  \mathbb{E}[X] = \sum_{i} x_i \, p_i.
\]

The single most useful tool in the analysis of randomized algorithms is
\emph{linearity of expectation}: for random variables $X_1, X_2, \ldots,
X_n$ (which need not be independent),
\begin{equation*}
  \mathbb{E}\!\Bigl[\sum_{i=1}^{n} X_i\Bigr]
  = \sum_{i=1}^{n} \mathbb{E}[X_i].
  \tag{1.6}
\end{equation*}

\begin{example}[The Sailor Problem]
A ship arrives at a port, and the 40 sailors on board go ashore for
revelry. Later at night, the 40 sailors return to the ship and, in their
state of inebriation, each chooses a random cabin to sleep in. What is the
expected number of sailors sleeping in their own cabins?

Let $X_i$ be 1 if the $i$th sailor chooses her own cabin, and 0 otherwise.
We wish to compute $\mathbb{E}[\sum_{i=1}^{40} X_i]$. By linearity of
expectation, this is $\sum_{i=1}^{40} \mathbb{E}[X_i]$. Since the cabins
are chosen at random, $\mathbb{E}[X_i]=1/40$. Thus the expected number is
$\sum_{i=1}^{40} 1/40 = 1$.
\end{example}

\subsection{Variance and deviations}

The \emph{variance} of a random variable $X$ measures its spread:
\[
  \mathrm{Var}(X) = \mathbb{E}[(X - \mathbb{E}[X])^2]
                  = \mathbb{E}[X^2] - (\mathbb{E}[X])^2.
\]
The \emph{standard deviation} is $\sigma(X) = \sqrt{\mathrm{Var}(X)}$.

\emph{Chebyshev's inequality} bounds the probability that $X$ deviates
significantly from its mean:
\begin{equation*}
  \Pr\bigl[|X - \mathbb{E}[X]| \ge t\bigr]
  \le \frac{\mathrm{Var}(X)}{t^2}.
  \tag{1.7}
\end{equation*}

This inequality will be applied frequently in Chapter~4, where we develop
sharper tail bounds for sums of independent random variables.

\subsection{The union bound}

A simple yet powerful tool is the \emph{union bound} (or \emph{Boole's
inequality}): for events $\mathcal{E}_1, \mathcal{E}_2, \ldots,
\mathcal{E}_n$ (not necessarily independent),
\begin{equation*}
  \Pr\!\Bigl[\bigcup_{i=1}^{n} \mathcal{E}_i\Bigr]
  \le \sum_{i=1}^{n} \Pr[\mathcal{E}_i].
  \tag{1.8}
\end{equation*}
We used the union bound implicitly in the analysis of Karger's algorithm
when bounding the probability that any edge of the min-cut is contracted.

%==========================================================================
\section{The Randomized Approach}
\label{sec:randomized-approach}

A randomized algorithm is an algorithm that can toss coins during its
execution. Formally, we model a randomized algorithm $R$ as a deterministic
algorithm that takes as an additional input a sequence
$r = r_1, r_2, \ldots$ of independent, uniformly distributed random
bits. The output and running time of $R$ are therefore random variables
that depend on $r$.

Randomness serves several distinct roles in algorithm design. We summarize
the four principal motivations before examining each in turn.

\begin{enumerate}
\item \textbf{Improved average-case complexity.} Even when efficient
deterministic algorithms exist, randomization can yield better expected
running time or simpler analysis.
\item \textbf{No efficient deterministic algorithm known.} For certain
problems, randomization provides the only known polynomial-time solution.
\item \textbf{Simplicity.} Randomized algorithms are often dramatically
easier to design, implement, and analyze than their deterministic
counterparts.
\item \textbf{Adversarial immunity.} Randomness shields an algorithm
against worst-case inputs crafted by an adversary who knows the algorithm's
strategy.
\end{enumerate}

\subsection{Improved average-case complexity}

The classical example is sorting. Deterministic comparison-based sorting
achieves $O(n \log n)$ worst-case time, but the algorithms that achieve
this bound (merge sort, heap sort) have overheads that make them slower in
practice than quicksort. Deterministic quicksort, however, has $O(n^2)$
worst-case time. Randomized quicksort resolves this dilemma: it always
produces a correctly sorted array (Las Vegas), yet its expected running
time is $O(n \log n)$, and its practical performance is excellent because
the constant factors are small. The random choice of pivot eliminates the
pathological inputs that defeat deterministic quicksort.

A second example is randomized selection (Section~4.3). The deterministic
median-finding algorithm of Blum et al.\ runs in $O(n)$ time but has large
constants. A simple randomized algorithm achieves $O(n)$ expected time
with much smaller constants by choosing pivots at random.

Further examples include:
\begin{itemize}
\item \textbf{Hash tables.} Universal hashing (Section~9.3) guarantees
$O(1)$ expected time for insertions, deletions, and lookups, even under
adversarial keys. Deterministic hash tables require strong assumptions
about the input distribution or more complex data structures.
\item \textbf{Skip lists.} Skip lists (Section~9.2) achieve $O(\log n)$
expected time for all dictionary operations with a simple randomized
data structure that competes with balanced trees.
\item \textbf{Randomized incremental algorithms.} Building a convex hull
of $n$ points (Section~10.2) in expected $O(n \log n)$ time using
randomized incremental insertion avoids the complex case analysis of
deterministic algorithms like Graham scan.
\item \textbf{Load balancing.} Throwing $m$ balls into $n$ bins uniformly
at random (Chapter~4) yields a maximum load of $O(m/n + \sqrt{m \log n})$.
When $m \approx n$, the maximum load is $O(\log n / \log \log n)$ with
high probability, far smaller than the deterministic worst case.
\end{itemize}

\subsection{No efficient deterministic algorithm known}

Before 2002, no polynomial-time deterministic algorithm was known for
primality testing. The randomized Miller--Rabin test (Chapter~15) runs in
$O(k \log^3 n)$ time with error probability at most $4^{-k}$, making it the
method of choice for decades. The deterministic AKS algorithm, discovered
in 2002, runs in $O(\log^{7.5} n)$ time but is significantly slower in
practice.

Further examples include:
\begin{itemize}
\item \textbf{Polynomial identity testing.} The Schwartz--Zippel lemma
(Section~8.2) provides a randomized algorithm for testing whether a
polynomial is identically zero that is dramatically faster than any
known deterministic method.
\item \textbf{Interactive proof systems.} The class $\mathbf{IP}$ equals
$\mathbf{PSPACE}$ only when randomness and interaction are both allowed
(Section~8.8). Without randomness, interactive proofs collapse to
$\mathbf{NP}$.
\item \textbf{Approximate counting.} Counting the number of satisfying
assignments to a DNF formula (Section~12.2) admits a fully polynomial
randomized approximation scheme, while no deterministic approximation
scheme with comparable guarantees is known.
\item \textbf{Perfect matching in parallel.} Finding a perfect matching
in a graph can be done in $\mathbf{RNC}$ (randomized NC) using the
Isolating Lemma (Section~13.4). No $\mathbf{NC}$ algorithm (deterministic
polylogarithmic-time parallel) is known.
\end{itemize}

\subsection{Simplicity}

Karger's min-cut algorithm (Section~1.4.2) is a striking illustration. The
deterministic algorithms for min-cut are based on network flows and require
sophisticated data structures. Karger's randomized algorithm, by contrast,
is a dozen lines of code: repeatedly contract a random edge. The analysis
is elegant and the implementation is trivial.

Further examples include:
\begin{itemize}
\item \textbf{Freivalds' algorithm.} Verifying the product of two $n
\times n$ matrices (Section~8.1) can be done in $O(n^2)$ time with a
randomized algorithm that multiplies a vector by each matrix. The
deterministic $O(n^2)$ verification algorithm requires a complicated
reduction; the randomized version is a few lines of code.
\item \textbf{Luby's maximal independent set.} Luby's algorithm
(Section~13.3) for finding a maximal independent set in a graph is a
simple randomized parallel algorithm that runs in $O(\log n)$ rounds.
Deterministic parallel algorithms for MIS are significantly more complex.
\item \textbf{Randomized incremental construction.} Deterministic
algorithms for problems like convex hull, Delaunay triangulation, and
linear programming are often complex, requiring careful case analysis.
The randomized incremental approach (Chapter~10) replaces this complexity
with a simple random sampling strategy.
\item \textbf{Tree data structures.} Random treaps and skip lists
(Sections~8.1--8.2) provide simple alternatives to balanced binary search
trees, replacing rotation invariants with randomization.
\end{itemize}

\subsection{Adversarial immunity}

An adversary who knows a deterministic algorithm's code can construct a
worst-case input tailored to exploit its predictable behavior. Randomness
removes this vulnerability: the adversary cannot predict the algorithm's
random choices.

Further examples include:
\begin{itemize}
\item \textbf{Online paging.} No deterministic online algorithm can achieve
a competitive ratio better than $k$ for $k$-page caching. The randomized
marking algorithm (Section~14.1) achieves $O(\log k)$ competitive ratio
against an oblivious adversary.
\item \textbf{Load balancing.} When jobs are assigned to machines by random
choice, the maximum load on any machine is concentrated near the average,
regardless of the job sizes---no adversary can force a bad assignment.
\item \textbf{Byzantine agreement.} Deterministic Byzantine agreement
requires $n \ge 3t+1$ processors to tolerate $t$ faults. Randomized
Byzantine agreement (Section~13.6) can tolerate $t < n/3$ faults with
expected constant-round protocols, circumventing the deterministic lower
bound.
\item \textbf{Quicksort.} Deterministic quicksort is vulnerable to an
adversary who supplies a pre-sorted array, causing $O(n^2)$ behavior.
Randomized quicksort eliminates this attack by choosing pivots randomly.
\item \textbf{Routing on the hypercube.} The randomized oblivious routing
algorithm for the hypercube (Section~5.2) achieves $O(\log n)$ expected
delay for any permutation, while deterministic oblivious routing can
require $\Omega(\sqrt{n})$ delay for some permutations.
\end{itemize}

\subsection{Las Vegas and Monte Carlo algorithms}

Randomized algorithms differ in the nature of their randomness.

\begin{definition}[Las Vegas Algorithm]
A \emph{Las Vegas algorithm} is a randomized algorithm that always produces
the correct output. Its running time is a random variable, and we
characterize it by its expected value.
\end{definition}

\begin{definition}[Monte Carlo Algorithm]
A \emph{Monte Carlo algorithm} is a randomized algorithm that may produce an
incorrect output with some small probability. Its running time is typically
a deterministic bound (e.g., polynomial in the input size).
\end{definition}

Monte Carlo algorithms for decision problems (problems with YES/NO answers)
come in two flavors:
\begin{itemize}
\item \textbf{One-sided error:} The algorithm may err only when it outputs
NO (or only when it outputs YES), but never in the other case.
\item \textbf{Two-sided error:} The algorithm may err regardless of the
correct answer.
\end{itemize}

A Las Vegas algorithm is a special case of a Monte Carlo algorithm with
error probability zero. Conversely, whenever solutions can be verified
efficiently, a Monte Carlo algorithm can be converted into a Las Vegas
algorithm by repeating and checking (see Exercise~1.2 below).

\subsection{The min-cut algorithm: a first example}

We now illustrate the power of randomization with a beautiful algorithm for
the min-cut problem. Let $G$ be a connected, undirected multigraph with $n$
vertices. A \emph{cut} in $G$ is a set of edges whose removal disconnects
$G$. A \emph{min-cut} is a cut of minimum cardinality.

The following randomized algorithm, due to Karger, is disarmingly simple.
We repeat the following step: pick an edge uniformly at random and merge the
two vertices at its endpoints. If there are several edges between some pairs
of newly formed vertices, retain them all. Edges between vertices that are
merged are removed, so that there are never any self-loops. We refer to this
process of merging the two endpoints of an edge into a single vertex as the
\emph{contraction} of that edge. With each contraction, the number of
vertices decreases by one. The crucial observation is that an edge
contraction does not reduce the min-cut size: every cut in the graph at any
intermediate stage is a cut in the original graph. The algorithm continues
the contraction process until only two vertices remain; the set of edges
between these two vertices is output as a candidate min-cut.

\begin{figure}[htbp]
\centering
\begin{tikzpicture}[scale=0.8,
  every node/.style={circle,draw,minimum size=8mm}]
  \node (1) at (0,1)   {1};
  \node (2) at (1.5,1) {2};
  \node (3) at (0.75,2){3};
  \node (4) at (0,0)   {4};
  \node (5) at (1.5,0) {5};
  \draw (1)--(2); \draw (1)--(3); \draw (1)--(4);
  \draw (2)--(3); \draw (2)--(5); \draw (3)--(4);
  \draw (3)--(5); \draw (4)--(5);
  \node at (0.75,-0.8){\small (a) Original graph};

  \node (12) at (4,1)   {$\{1,2\}$};
  \node (3b) at (4.75,2){3};
  \node (4b) at (4,0)   {4};
  \node (5b) at (4.75,0){5};
  \draw (12)--(3b); \draw (12)--(4b); \draw (12)--(5b);
  \draw (3b)--(4b); \draw (3b)--(5b); \draw (4b)--(5b);
  \node at (4.375,-0.8){\small (b) After contracting $(1,2)$};
\end{tikzpicture}
\caption{A step in the min-cut algorithm; the effect of contracting
  edge $e=(1,2)$ is shown.}
\label{fig:contraction}
\end{figure}

Does this algorithm always find a min-cut? No---it may err. But we can
bound the probability of error. Let $k$ be the min-cut size. We fix our
attention on a particular min-cut $C$ with $k$ edges. Clearly $G$ has at
least $kn/2$ edges; otherwise there would be a vertex of degree less than
$k$, and its incident edges would be a min-cut of size less than $k$.

Let $e_i$ denote the event of not picking an edge of $C$ at the $i$th step,
for $1 \le i \le n-2$. The probability that the edge randomly chosen in the
first step is in $C$ is at most $k/(nk/2)=2/n$, so that $\Pr[e_1] \ge
1-2/n$. Assuming that $e_1$ occurs, during the second step there are at
least $k(n{-}1)/2$ edges, so the probability of picking an edge in $C$ is at
most $2/(n{-}1)$, so that $\Pr[e_2 \mid e_1] \ge 1-2/(n{-}1)$. At the $i$th
step, the number of remaining vertices is $n-i+1$. The min-cut size is still
at least~$k$, so the graph has at least $k(n{-}i{+}1)/2$ edges. Thus,
$\Pr[e_i \mid \bigcap_{j=1}^{i-1} e_j] \ge 1-2/(n{-}i{+}1)$.

The probability that no edge of $C$ is ever picked is
\begin{equation*}
  \Pr\!\Bigl[\bigcap_{i=1}^{n-2} e_i\Bigr]
  = \prod_{i=1}^{n-2}\Bigl(1-\frac{2}{n-i+1}\Bigr)
  = \frac{2}{n(n-1)}.
  \tag{1.1}
\end{equation*}

The probability of discovering a particular min-cut is therefore at least
$2/n^2$. To reduce the failure probability, we repeat the algorithm
$n^2/2$ times, making independent random choices each time. The probability
that a min-cut is not found in any attempt is at most
\begin{equation*}
  \Bigl(1-\frac{2}{n^2}\Bigr)^{n^2/2} < \frac{1}{e}.
  \tag{1.2}
\end{equation*}

Further repetitions make the failure probability arbitrarily small---the
only consideration being that repetitions increase running time. Note the
extreme simplicity of this randomized algorithm compared to the
deterministic algorithms for min-cut, which are based on network flows and
are considerably more complex.

\begin{exercise}
Suppose that at each step of the min-cut algorithm, instead of choosing a
random edge for contraction we choose two vertices at random and coalesce
them into a single vertex. Show that there are inputs on which the
probability that this modified algorithm finds a min-cut is exponentially
small.
\end{exercise}

%% ----- C++ : Min-Cut -----
\subsection{C++ Implementation}

The following implementation demonstrates Karger's algorithm.  A multigraph
is represented with weighted adjacency lists; the weight of edge $(u,v)$
equals its multiplicity.

\begin{lstlisting}[
  caption={Graph data structure for the min-cut algorithm.},
  label={lst:graph}]
#include <vector>
#include <unordered_map>
#include <algorithm>
#include <numeric>
#include <climits>
#include <random>
#include <chrono>
#include <iostream>
#include <cmath>

class RandomEngine {
public:
  static RandomEngine& instance() {
    static RandomEngine e; return e;
  }
  int uniform_int(int lo, int hi) {
    return std::uniform_int_distribution<int>(lo,hi)(eng_);
  }
  double uniform_real(double lo=0,double hi=1) {
    return std::uniform_real_distribution<double>(lo,hi)(eng_);
  }
private:
  RandomEngine()
    : eng_(std::chrono::steady_clock::now()
           .time_since_epoch().count()) {}
  std::mt19937 eng_;
};

struct Graph {
  int n;
  std::unordered_map<int,
    std::unordered_map<int,int>> adj;
  explicit Graph(int n) : n(n) {}
  void add_edge(int u, int v, int c=1) {
    adj[u][v]+=c; adj[v][u]+=c;
  }
  int edge_count() const {
    int t=0;
    for (auto& [u,nb] : adj)
      for (auto& [v,w] : nb)
        if (u<v) t+=w;
    return t;
  }
};
\end{lstlisting}

\begin{lstlisting}[
  caption={Karger's min-cut algorithm.},
  label={lst:karger}]
void contract(auto& adj, int u, int v) {
  for (auto& [w,wt] : adj[v])
    if (w!=u) { adj[u][w]+=wt; adj[w][u]+=wt; }
  adj.erase(v);
  for (auto& [_,nb] : adj) nb.erase(v);
  adj[u].erase(u);            // remove self-loops
}

int karger_min_cut(Graph& g) {
  auto adj = g.adj;
  std::vector<int> vs(g.n);
  std::iota(vs.begin(), vs.end(), 0);
  while (vs.size()>2) {
    // collect all edges (with multiplicity)
    std::vector<std::pair<int,int>> E;
    for (auto& [u,nb] : adj)
      for (auto& [v,w] : nb)
        if (u<v)
          for (int i=0;i<w;i++) E.push_back({u,v});
    if (E.empty()) break;
    auto [u,v] = E[RandomEngine::instance()
                        .uniform_int(0,E.size()-1)];
    contract(adj,u,v);
    vs.erase(std::remove(vs.begin(),vs.end(),v),
             vs.end());
  }
  return vs.size()==2 ? adj[vs[0]][vs[1]] : 0;
}

int karger_repeated(Graph& g, int T) {
  int best = INT_MAX;
  for (int t=0;t<T;t++)
    best = std::min(best, karger_min_cut(g));
  return best;
}
\end{lstlisting}

\begin{lstlisting}[
  caption={Driver: empirical verification of the theoretical bounds.},
  label={lst:mincut_demo}]
int main() {
  Graph g(5);
  g.add_edge(0,1); g.add_edge(0,2); g.add_edge(0,3);
  g.add_edge(1,2); g.add_edge(1,3); g.add_edge(2,3);
  g.add_edge(2,4); g.add_edge(3,4);

  int n=g.n, T=(n*n)/2;
  std::cout<<"Single trial : "<<karger_min_cut(g)<<"\n";
  std::cout<<T<<" trials     : "<<karger_repeated(g,T)<<"\n";

  double p=2.0/(n*(n-1));
  std::cout<<"P(success) >= 2/n(n-1) = "<<p<<"\n";
  std::cout<<"P(failure in "<<T<<" trials) <= "
           <<std::pow(1-p,T)<<"\n";
  std::cout<<"1/e = "<<1.0/std::exp(1.0)<<"\n";
}
\end{lstlisting}

\paragraph{Complexity.}
Each trial of the algorithm runs in $O(n^2)$ time with the adjacency-map
representation.  With $n^2/2$ trials the total expected running time is
$O(n^4)$.

%% ----- C++ : Las Vegas -----
\subsection{C++ Implementation: Randomized QuickSort}

Randomized QuickSort is the canonical Las Vegas algorithm: it always
produces a correctly sorted array, yet its running time is a random variable
with $\mathbb{E}[T(n)]=O(n\log n)$.

\begin{lstlisting}[
  caption={Randomized QuickSort---a Las Vegas algorithm.},
  label={lst:quicksort}]
int partition(std::vector<int>& a, int lo, int hi) {
  int pi = RandomEngine::instance().uniform_int(lo,hi);
  std::swap(a[pi], a[hi]);
  int pivot = a[hi], i = lo-1;
  for (int j=lo; j<hi; j++)
    if (a[j]<=pivot) std::swap(a[++i], a[j]);
  std::swap(a[i+1], a[hi]);
  return i+1;
}

void rand_quicksort(std::vector<int>& a,int lo,int hi){
  if (lo<hi) {
    int p = partition(a,lo,hi);
    rand_quicksort(a,lo,p-1);
    rand_quicksort(a,p+1,hi);
  }
}
\end{lstlisting}

%% ----- C++ : Monte Carlo -----
\subsection{C++ Implementation: Monte Carlo---Estimating \texorpdfstring{$\pi$}{pi}}

The following Monte Carlo algorithm estimates $\pi$ by sampling.  It has
two-sided error: the estimate may be too high or too low, but the error
decreases as $O(1/\sqrt{n})$.

\begin{lstlisting}[
  caption={Monte Carlo estimation of $\pi$ (two-sided error).},
  label={lst:pi}]
double estimate_pi(int n) {
  int inside = 0;
  for (int i=0;i<n;i++) {
    double x = RandomEngine::instance()
                   .uniform_real(-1,1);
    double y = RandomEngine::instance()
                   .uniform_real(-1,1);
    if (x*x+y*y<=1.0) inside++;
  }
  return 4.0*inside/n;
}
\end{lstlisting}

\begin{table}[htbp]
\centering
\begin{tabular}{rrr}
\toprule
\textbf{Samples} & \textbf{Estimate} & \textbf{Error} \\
\midrule
    100 & 3.12000 & 0.02159 \\
  1\,000 & 3.14400 & 0.00241 \\
  10\,000 & 3.13840 & 0.00319 \\
 100\,000 & 3.14076 & 0.00083 \\
  \bottomrule
\end{tabular}
\caption{Typical output of the $\pi$ estimator.}
\end{table}

%% ----- C++ : MC $\to$ LV -----
\subsection{Converting Between Las Vegas and Monte Carlo}

A Las Vegas algorithm can be converted to a Monte Carlo algorithm by
truncation.  If the Las Vegas algorithm has expected running time
$T(n)$, run it for $2T(n)$ steps and output an arbitrary default
answer if it has not finished.  By Markov's inequality, the
probability of premature truncation is at most~$1/2$, so the resulting
Monte Carlo algorithm errs with probability at most~$1/2$ and always
runs in $O(T(n))$ time.  Conversely, when solutions can be
\emph{verified} efficiently, a Monte Carlo algorithm can be turned
into a Las Vegas algorithm by repeating until verification succeeds.

\begin{lstlisting}[
  caption={Monte Carlo $\to$ Las Vegas via verification.},
  label={lst:conversion}]
// Monte Carlo: correct w.p.~2/3
std::pair<int,bool> mc(int x) {
  bool ok = RandomEngine::instance()
                 .uniform_real() < 2.0/3.0;
  return ok ? std::make_pair(2*x,true)
            : std::make_pair(3*x,false);
}

bool verify(int x, int r) { return r==2*x; }

// Las Vegas wrapper
std::pair<int,int> mc_to_lv(int x) {
  int tries=0;
  while (true) {
    tries++;
    auto [r,ok] = mc(x);
    if (verify(x,r)) return {r,tries};
  }
}
\end{lstlisting}

\begin{remark}
If the single-trial success probability is~$p$, the number of independent
trials until the first success follows a \emph{geometric distribution} with
mean~$1/p$.  Hence the expected running time of the Las Vegas wrapper is
$(T(n)+t(n))/p$.
\end{remark}

%==========================================================================

\end{document}
